% SPDX-License-Identifier: CC-BY-4.0
% Copyright (C) 2026 Fabian Januszewski, Abhijit (the Langlandscape team)
% Report section: what the Lean formalization assumes. \input-able; needs report/hyp-preamble.tex and
% report/hypotheses.bib.  Sources: lean/Langlandscape/HYPOTHESES.md, SORRIES.md, Langlandscape/Hyp/*.lean,
% Audit/PrintAxioms.lean, lean/Langlandscape/audits/*.md, mathematics/audits/*.md, mathematics/chapters/s0*,
% plan/implementation/*.md; snapshot origin/wip-implementation d214fb58 plus impl/s04-int15 cd0a268 (2026-11-15);
% stage-6 rows and counts refreshed at wip-implementation 99f5500f (2026-11-15, S06-M10b).
% Refresh procedure: section sec:hyp-refresh.

%======================================================================================================
\section{How hypotheses are handled}\label{sec:hyp-intro}

\paragraph{Snapshot.}
This section describes the Lean project \lname{lean/Langlandscape/} at the end of stage~4
(state of 2026-11-15, \emph{provisional}), with the stage-6 hypotheses added from \lname{wip-implementation} at
\lname{99f5500f} (2026-11-15; registry rows of the S06 registry worker, Lean audit S06 part 2). Two trees are used:
the integration branch \lname{wip-implementation} at \lname{d214fb58}, and the stage-4 tree with the final audit
fixes, \lname{impl/s04-int15} at \lname{cd0a268}. The second tree is not yet merged; it differs from the first, as far as
hypotheses go, only in \lname{Hyp.InvLimExact} (a connectivity premise was added, S04-L17c, audit R13), in a scope
note on \lname{Hyp.MordLambdaFree} (R14), and in eighteen \lname{\_uncond} theorems that remove
\lname{Hyp.PadicOrdDecomp} from the tower results (R11). Everything stated below holds for \lname{cd0a268} unless
marked otherwise. Section~\ref{sec:hyp-refresh} says how to refresh the document after the stage-4 final gate.

\paragraph{Hypotheses are explicit arguments.}
A mathematical input that is used but not proved in Lean is a \emph{registered hypothesis}: a definition
\begin{center}\lname{def Langlandscape.Hyp.X (parameters) : Prop := \dots}\end{center}
in the file \lname{Langlandscape/Hyp/X.lean}. A theorem that needs it takes a proof of \lname{Hyp.X \dots} as an
ordinary, explicit argument, for example
\begin{center}
\lname{theorem ribet\_exact\_of\_SA\_HS (R) (hSA : Hyp.SA d q hq) (hHS : Hyp.HS d) (hBT : Hyp.BTtree q) (hU : \dots) : \dots}
\end{center}
Hypotheses are never Lean \lname{axiom}s and never global type-class instances (\lname{plan/implementation/OVERVIEW.md}
\S3.6). Hypotheses of the imported trace-formula project are reused under their own name and not redefined
(\S\ref{sec:hyp-import}). The registry of all hypotheses is \lname{lean/Langlandscape/HYPOTHESES.md}, which only the
stage master edits. The registry of \lname{sorry}s, \lname{lean/Langlandscape/SORRIES.md}, is empty: there is no
\lname{sorry} in \lname{lean/Langlandscape/} (S01-X1a onwards; \emph{checked} by every gate).

Besides the registered hypotheses there are \emph{standing hypotheses}: arithmetic side conditions bundled in
\lname{Prop}-valued structures, e.g.\ \lname{Ordinary.TowerHyp d M p} ($p\nmid D$, $M\neq 0$, $p\nmid M$,
$(Mp^s,D)=1$) and \lname{Ordinary.RibetHyp d q M p} (adds $q\nmid D$, $q\nmid Mp^s$), and premises such as
\lname{hconn : Tree.BGConnected \dots}. These are conditions on the input data of a theorem, like ``$p$ is prime'';
they are not mathematical results and are not registered.

\paragraph{The axiom audit.}
\lname{lean/Langlandscape/Audit/PrintAxioms.lean} runs \lname{\#print axioms} on every headline declaration
(hand-maintained list; candidates generated by \lname{tools/gen\_axioms.py}, counted by \lname{tools/axcount.sh}).
\emph{Computed} results of the gates:
\begin{center}\small
\begin{tabular}{@{}llrrl@{}}\toprule
gate & commit & printed/expected & non-standard & \lname{sorryAx}\\\midrule
TITAN \lname{s03-final-1} & \lname{5bda8569} & 519/519 & 0 & 0\\
TITAN \lname{s04-gate1-1} & \lname{a8cbd708} & 558/558 & 0 & 0\\
CONCEPTD \lname{s04-int15-1} & \lname{cd0a268} & 592/592 & 0 & 0\\\bottomrule
\end{tabular}
\end{center}
At \lname{cd0a268}: 580 declarations use exactly \lname{propext}, \lname{Classical.choice}, \lname{Quot.sound};
10 use a strict subset (\lname{Hyp.HsiehCR} only \lname{propext}; \lname{Hyp.LambdaJL} and eight
\lname{Graph}/\lname{CommAlg} names \lname{[propext, Quot.sound]}); 2 (\lname{Graph.ofBip}, \lname{Graph.collapse2})
use no axiom. Logs: \lname{lean/Langlandscape/.titan/out/axcount.txt}, \lname{pa.log} on
\lname{origin/conceptd/s04-int15-1}; \lname{c-axcount.log} on \lname{origin/titan/s04-gate1-1}.

\paragraph{What the two audits say together.}
\lname{\#print axioms} lists every axiom on which a declaration depends, transitively through all imported modules,
including Mathlib and the trace-formula import. A \lname{sorry} anywhere in that dependency cone would show up as
\lname{sorryAx}; a hidden assumption would show up as a non-standard axiom. Neither occurs. Registered hypotheses do
\emph{not} show up there, by design: they are arguments, not axioms. So the axiom audit certifies that nothing is
assumed \emph{implicitly}, and the type of each theorem shows exactly what is assumed \emph{explicitly}.

\begin{readbox}[How to read a formal theorem of this project]\label{box:hyp-read}
\small Suppose the Lean theorem \lname{T} has the type
\lname{(h$_1$ : Hyp.X$_1$ \dots) $\to$ \dots $\to$ (h$_k$ : Hyp.X$_k$ \dots) $\to$ (standing premises) $\to$ C}, and
the axiom audit for \lname{T} is clean. Then the following is \emph{formalized}:
\begin{enumerate}[leftmargin=1.4em]
\item \textbf{What is proved.} In Lean's type theory with \lname{propext}, \lname{Classical.choice} and
  \lname{Quot.sound}, the conclusion \lname{C} follows from the propositions \lname{Hyp.X$_i$} \emph{as defined in
  Lean} and the standing premises. Nothing else is used.
\item \textbf{What the reader must check.} That each \lname{Hyp.X$_i$}, read literally, says what the cited source
  proves (the entries of \S\ref{sec:hyp-entries} give the statement, the source, and the audit that checked
  fidelity), and that \lname{C} says what the chapter claims (the \lname{\textbackslash lean\{\dots\}} marks).
\item \textbf{Vacuity.} If some \lname{Hyp.X$_i$} were false, \lname{T} would be vacuously true. Each entry therefore
  records the non-vacuity evidence: a closed instance, a toy instance, or an argument why the statement is true.
\item \textbf{Schematic hypotheses} (\lname{JLClassical}, \lname{LambdaJL}, \lname{HidaControl}, \lname{HsiehCR})
  take abstract data, such as an abstract Hecke module standing in for $S_2(\Gamma_0(N))$. They mean the
  mathematical theorem \emph{only} when instantiated with the true objects, which cannot be done in Mathlib v4.35.
  No theorem may be advertised as ``under classical JL'' without this caveat (OVERVIEW \S8, F6).
\item \textbf{Status words.} \emph{formalized modulo \lname{Hyp.X}} = item 1 holds;
  \emph{formalized} = no registered hypothesis in the type (other explicit hypotheses of the statement, e.g.\ connectivity,
  remain); \emph{discharged} = a closed Lean proof of \lname{Hyp.X} exists (so every consumer has a form without that argument).
\end{enumerate}
\end{readbox}

%======================================================================================================
\section{Overview}\label{sec:hyp-overview}

\paragraph{Kinds.}
Each hypothesis is classified by \emph{why} it is assumed. The classification is ours (it is not a registry column);
it is recorded in \lname{mathematics/report/tools/hyp\_table.py}.
\begin{description}[leftmargin=1.2em,font=\normalfont]
\item[\kindlit] \emph{Known literature result}, cited with a precise reference; formalizing it needs Mathlib
  infrastructure that does not exist (adelic groups, Bass--Serre theory, local division algebras, class number formulas).
\item[\kindinfra] \emph{Proved on paper in this project} (chapter mark \stproved{} or \stpmc); the Lean proof needs
  infrastructure that is in reach but not yet built.
\item[\kindoos] \emph{Target statement, formalization out of scope}: the classical and $\Lambda$-adic JL
  correspondences, stated schematically and cited.
\item[\kindcond] \emph{Condition}: a property of the input data (a residual representation), not a theorem;
  it is never discharged, only checked case by case.
\item[\kinddis] \emph{Discharged}: a closed Lean proof exists.
\item[\kindplan] \emph{Planned} for stages 5--6 (\S\ref{sec:hyp-planned}); not yet in Lean.
\end{description}
No registered hypothesis is open or conjectural mathematics: every one is either a theorem in the literature, a
theorem proved in our chapters, a target statement that is a known theorem, or a condition.

\paragraph{Counts} (\emph{computed} from the registry at \lname{2323a252} by \lname{tools/hyp\_table.py}; method (a): rows \lname{| `Langlandscape.Hyp.} of \lname{HYPOTHESES.md}, counted by \lname{grep -c}, plus the reused \lname{GL2TF.Hyp.EichlerMass} row): 29 registered
hypotheses (28 in \lname{Langlandscape.Hyp}, plus \lname{GL2TF.Hyp.EichlerMass} reused from the import; the alias
\lname{Langlandscape.Hyp.EichlerMass} is not counted separately). Of these, 26 are in Table~\ref{tab:hyp-status} (17 from stages
1--4, 9 from stage~6); the 3 stage-5 hypotheses \lname{Hyp.X.char}, \lname{Hyp.X.gal}, \lname{Hyp.X.iso} are registered but still
described in \S\ref{sec:hyp-planned}. Table by kind: 13~\kindlit, 7~\kindinfra, 2~\kindoos, 1~\kindcond, 3~\kinddis. By status:
3 discharged (\lname{PadicOrdDecomp}; \lname{RamNrdSurj}; \lname{ResidualBurnside}, only given absolute irreducibility, see below), 23 open, of which 3
are marked schematic in the registry. After that registry state, \lname{Hyp.BrandtStabFree} was discharged by
\lname{Langlandscape.Hyp.brandtStabFree\_holds} (an external draft, mistral-large-4-0 via opencode on the second team member's
machine, outside our agent process; integrated into the library, independently audited (Fable, PASS) and gated by our agents). With it
the counts read: by kind 6~\kindinfra{} and 4~\kinddis; by status 4 discharged (three closed proofs and \lname{ResidualBurnside}
given absolute irreducibility), 22 open; 14 of the 26 table rows have instances. 13 of the 26 table rows (method (b): rows of \lname{generated/hyp-table.tex}) have instances in the registry's sense (\lname{EichlerMass}: two closed instances in the
import's test library; \lname{PadicOrdDecomp}, \lname{RamNrdSurj}: proved; \lname{HsiehCR} and the nine stage-6 rows \lname{HsiehGor},
\lname{OldDecomp}, \lname{HsiehGorFull}, \lname{FreeRankTwo}, \lname{ResidualBurnside}, \lname{TamSpec}, \lname{HidaControlGL2}, \lname{KimOta}, \lname{KimOtaRT}: toy instances in \lname{LanglandscapeTests/S06/Hyp.lean}). Toy instances are consistency (non-vacuity) instances on abstract data
only, \emph{not} evidence for the mathematical statement.

\paragraph{What is assumed in stage 6.} \lname{Hyp.FreeRankTwo} is a Galois-theoretic input in hypothesis form: given \lname{Hyp.ResidualBurnside}, \lname{Hyp.FreeRankTwo} is equivalent to the inclusion $\tq_q\subseteq\cc_q$; both directions are proved in Lean (\lname{Defect.galois}, \lname{Defect.freeRankTwo\_of\_le}, packaged as \lname{Defect.freeRankTwo\_iff\_le}). The Galois half is therefore not formalized independently: its arithmetic content (that the old-form representation $\rho_M$ exists at level $M$, steps 1--3 of \lname{prop:s06-galois}) is the cited hypothesis (Lean audit S06 part 2, R1). So the Lean formalizes only the trace-form step (\lname{Defect.galois} via
\lname{CommAlg.eq\_one\_of\_trace\_mul\_eq}); the proposition itself is not a Lean theorem. Its kind \kindinfra{} in the table means
that the chapter proves the inclusion on paper (\stpmc). \lname{Hyp.ResidualBurnside} is discharged only given absolute irreducibility
of the residual representation ((H2); \lname{residualBurnside\_of\_absIrred}). Any claim that \lname{Defect.defect\_formula} is
formalized must read ``formalized modulo registered hypotheses, including the Galois half (\lname{Hyp.FreeRankTwo})''.
\lname{Hyp.KimOtaRT} is the inclusion of \lname{prop:s06-levM} under (Lev) of \lname{hyp:s06-H}: $M=1$, or $M$ squarefree with
$\bar\rho$ unramified at every $\ell\mid M$ (for $\ell\equiv-1\bmod p$ in our reading of \cite[Thm.~3.14]{KO}); $N(\bar\rho)>1$ is open.

\begin{table}[ht]\centering\small
\caption{Status of the registered hypotheses. Generated from \lname{HYPOTHESES.md} by
\lname{tools/hyp\_table.py}; ``instances'' = a closed proof or (toy) instance of the proposition exists.}
\label{tab:hyp-status}
\begin{tabular}{@{}lllll@{}}\toprule
Lean name & stage & kind & instances & registry status\\\midrule
\hyptablerows
\bottomrule\end{tabular}
\end{table}

\begin{figure}[!p]\centering
\resizebox{\textwidth}{!}{% SPDX-License-Identifier: CC-BY-4.0
% Copyright (C) 2026 Fabian Januszewski, Abhijit (the Langlandscape team)
% Dependency graph: registered hypotheses -> consuming Lean theorems -> stage-level statements.
% Status encoded twice: Okabe--Ito colour AND border style (survives greyscale printing).
\begin{tikzpicture}[
    x=1cm,y=1cm,font=\sffamily\scriptsize,
    every node/.style={align=center},
    hyp/.style={rectangle,rounded corners=1.5pt,minimum width=3.3cm,minimum height=0.46cm,inner sep=1.5pt,line width=0.9pt},
    lit/.style={hyp,draw=oiBlue,fill=oiBlue!14},
    dis/.style={hyp,draw=oiGreen,fill=oiGreen!18,double,double distance=0.9pt,line width=0.6pt},
    inf/.style={hyp,draw=oiOrange!85!black,fill=oiOrange!18,densely dashed,line width=1.1pt},
    oos/.style={hyp,draw=oiVerm,fill=oiVerm!14,line width=1.6pt},
    con/.style={hyp,draw=oiPurple!85!black,fill=oiPurple!16,densely dotted,line width=1.3pt},
    pln/.style={hyp,draw=oiSky!75!black,fill=oiSky!10,dashed,line width=0.9pt},
    thm/.style={rectangle,draw=black!70,fill=black!4,text width=4.3cm,minimum height=0.46cm,inner sep=2pt,line width=0.5pt,font=\sffamily\fontsize{6}{7}\selectfont},
    plt/.style={thm,dashed,draw=oiSky!75!black,fill=oiSky!6},
    fin/.style={rectangle,draw=black,fill=white,minimum width=3.4cm,inner sep=3pt,line width=1pt,text width=3.2cm},
    pfin/.style={fin,dashed,draw=oiSky!75!black},
    use/.style={-{Stealth[length=1.6mm]},line width=0.5pt,draw=black!65},
    puse/.style={use,dashed,draw=oiSky!75!black},
    intd/.style={use,dotted,draw=black!50},
    stg/.style={font=\sffamily\bfseries\tiny,text=black!60,anchor=east},
  ]
  % ------------------------------------------------ column 1: hypotheses
  \def\hx{0}
  \node[lit] (EM)  at (\hx,0)     {GL2TF.Hyp.EichlerMass};
  \node[oos] (JLC) at (\hx,-0.6)  {Hyp.JLClassical};
  \node[oos] (LJL) at (\hx,-1.2)  {Hyp.LambdaJL};
  \node[dis] (RN)  at (\hx,-2.0)  {Hyp.RamNrdSurj};
  \node[lit] (HS)  at (\hx,-2.6)  {Hyp.HS};
  \node[lit] (SA)  at (\hx,-3.2)  {Hyp.SA};
  \node[lit] (BT)  at (\hx,-3.8)  {Hyp.BTtree};
  \node[lit] (SC)  at (\hx,-4.4)  {Hyp.SerreCoh};
  \node[lit] (SF)  at (\hx,-5.0)  {Hyp.SerreFree};
  \node[lit] (EC)  at (\hx,-5.8)  {Hyp.EllipticCount};
  \node[inf] (BSF) at (\hx,-6.4)  {Hyp.BrandtStabFree};
  \node[inf] (IH)  at (\hx,-7.0)  {Hyp.IharaH0};
  \node[dis] (POD) at (\hx,-7.8)  {Hyp.PadicOrdDecomp};
  \node[inf] (ILE) at (\hx,-8.4)  {Hyp.InvLimExact};
  \node[inf] (MLF) at (\hx,-9.0)  {Hyp.MordLambdaFree};
  \node[lit] (HC)  at (\hx,-9.6)  {Hyp.HidaControl};
  \node[con] (CR)  at (\hx,-10.2) {Hyp.HsiehCR};
  \node[pln] (XC)  at (\hx,-11.0) {Hyp.X.char / X.gal / X.iso};
  \node[pln] (HG)  at (\hx,-11.6) {Hyp.HsiehGor};
  \node[pln,text width=3.6cm] (S6) at (\hx,-12.75)
     {OldDecomp, HsiehGorFull, FreeRankTwo, ResidualBurnside, TamSpec, HidaControlGL2, KimOta(RT)};
  \foreach \n/\y in {S01/0,S02/-2.0,S03/-5.8,S04/-7.8,S05/-11.0,S06/-11.6}{\node[stg] at (-1.75,\y) {\n};}

  % ------------------------------------------------ column 2: consuming theorems (Lean names, short)
  \def\tx{5.25}
  \node[thm] (tMass) at (\tx,0)     {Quat.card\_XB\_Utower,\\ Ribet.sum\_inv\_wx\_Ohat,\\ Ordinary.vp\_card\_XB\_Utower\_ge};
  \node[thm] (tMaz)  at (\tx,-1.05) {Ribet.mazur\_check};
  \node[thm] (tNrd)  at (\tx,-2.0)  {Quat.nrdSurj\_*, Tree.isQAdmissible\_*};
  \node[thm] (tBG)   at (\tx,-3.0)  {Tree.BGConnected\_of\_SA\_HS,\\ Tree.quotVEquiv/quotEEquiv\_of\_SA\_HS};
  \node[thm] (tRib)  at (\tx,-4.2)  {Ribet.ribet\_exact\_of\_SA\_HS,\\ ribet\_rank\_/ihara\_*\_of\_SA\_HS};
  \node[thm] (tPhi)  at (\tx,-6.4)  {Ribet.phiComb\_eisenstein\_Ohat/\_DR,\\ Ribet.brandtZ\_dvd};
  \node[thm] (tPOD)  at (\tx,-7.8)  {Ordinary.padicOrdDecomp\_holds\\ $\Rightarrow$ Ordinary.*\_uncond (tower)};
  \node[thm] (tLim)  at (\tx,-8.9)  {Ordinary.lim\_exact\_ribet(\_uncond)\\ (\texttt{hA}, \texttt{hconn} explicit)};
  \node[plt] (tXl)   at (\tx,-11.0) {S05: xlambda\_finite (planned)};
  \node[plt] (tDef)  at (\tx,-12.1) {S06: Setting.gor\_new, Defect.galois,\\ Defect.defect\_formula, cor\_fixed};

  % ------------------------------------------------ column 3: stage-level statements
  \def\fx{10.55}
  \node[fin] (F1) at (\fx,0.0)  {\textbf{S01} mass formula on the tower $U_s$ (formalized mod.\ EichlerMass)};
  \node[fin] (F3) at (\fx,-3.6) {\textbf{S03} Ribet exact sequence, Ihara, $\Phi$-Eisenstein, Mazur check at finite level};
  \node[fin] (F4) at (\fx,-8.3) {\textbf{S04} $\Lambda$-adic Ribet sequence ($\mathbb Z_p$-linear limit form)};
  \node[pfin] (F5) at (\fx,-11.0) {\textbf{S05} $\Lambda$-adic character group (planned)};
  \node[pfin] (F6) at (\fx,-12.4) {\textbf{S06} defect / Gorenstein formula (planned)};
  \node[fin,draw=oiVerm,line width=1.6pt] (FT) at (\fx,-1.4) {\textbf{Targets} classical and $\Lambda$-adic JL: stated, cited, no Lean consumer};

  % ------------------------------------------------ edges hyp -> theorem (Lean arguments)
  \draw[use] (EM.east) -- (tMass.west);
  \draw[use] (EM.east) -- (tMaz.west);
  \draw[use] (EC.east) -- (tMaz.west);
  \draw[use] (RN.east) -- (tNrd.west);
  \draw[use] (HS.east) -- (tBG.west);
  \draw[use] (SA.east) -- (tBG.west);
  \draw[use] (BT.east) -- (tBG.west);
  \draw[use] (HS.east) -- (tRib.west);
  \draw[use] (SA.east) -- (tRib.west);
  \draw[use] (BT.east) -- (tRib.west);
  \draw[use] (tNrd.south) -- node[right,font=\sffamily\tiny,inner sep=1pt]{IsQAdmissible} (tBG.north);
  \draw[use] (BSF.east) -- (tPhi.west);
  \draw[use] (POD.east) -- (tPOD.west);
  \draw[use] (POD.east) -- (tLim.west);
  \draw[use] (ILE.east) -- (tLim.west);
  % statement-only hypotheses (no library consumer): dotted, to the statement they belong to
  \draw[intd] (SC.east) to[out=0,in=200] (tRib.south west);
  \draw[intd] (SF.east) to[out=0,in=210] (tRib.south west);
  \draw[intd] (IH.east) to[out=0,in=-90] ($(tRib.south)+(0.6,0)$);
  \draw[intd] (MLF.east) to[out=0,in=180] (F4.west);
  \draw[intd] (HC.east) to[out=0,in=185] (F4.south west);
  \draw[intd] (JLC.east) to[out=0,in=180] (FT.west);
  \draw[intd] (LJL.east) to[out=10,in=190] (FT.west);
  % planned
  \draw[puse] (XC.east) -- (tXl.west);
  \draw[puse] (POD.south east) to[out=-20,in=160] (tXl.west);
  \draw[puse] (HG.east) -- (tDef.west);
  \draw[puse] (S6.east) -- (tDef.west);
  \draw[puse] (LJL.east) to[out=-15,in=150] (tDef.north west);
  \node[font=\sffamily\tiny,anchor=west,text=black!60] at ($(CR.east)+(0.1,0)$) {no consumer, none planned};
  % theorem -> stage statement
  \draw[use] (tMass.east) -- (F1.west);
  \draw[use] (tMaz.east) -- (F3.west);
  \draw[use] (tBG.east) -- (F3.west);
  \draw[use] (tRib.east) -- (F3.west);
  \draw[use] (tPhi.east) -- (F3.west);
  \draw[use] (tPOD.east) -- (F4.west);
  \draw[use] (tLim.east) -- (F4.west);
  \draw[puse] (tXl.east) -- (F5.west);
  \draw[puse] (tDef.east) -- (F6.west);
  \draw[puse] (F4.south) -- (F5.north);
  \draw[puse] (F5.south) -- (F6.north);

  % ------------------------------------------------ legend
  \begin{scope}[shift={(0,-14.75)},every node/.append style={minimum width=3.4cm}]
    \node[lit] at (0.0,0) {literature result (LIT)};
    \node[inf] at (4.4,0) {proved on paper (INF)};
    \node[oos] at (8.8,0) {target, out of scope (OOS)};
    \node[dis] at (0.0,-0.65) {discharged (DIS)};
    \node[con] at (4.4,-0.65) {condition (CON)};
    \node[pln] at (8.8,-0.65) {planned (PLN)};
  \end{scope}
  \begin{scope}[shift={(0,-16.1)}]
    \draw[use] (-1.6,0) -- ++(1,0) node[right,black]{explicit Lean argument};
    \draw[intd] (2.9,0) -- ++(1,0) node[right,black]{statement only (intended use)};
    \draw[puse] (8.0,0) -- ++(1,0) node[right,black]{planned (S05/S06)};
  \end{scope}
\end{tikzpicture}
}
\caption{Which hypotheses feed which theorems. Left: registered and planned hypotheses (colour and border style both
encode the kind). Middle: the Lean theorems that take them as arguments. Right: the stage-level statements. Solid
arrows are explicit Lean arguments at \lname{cd0a268}; dotted arrows join a statement-only hypothesis (no library
consumer) to the result it belongs to; dashed arrows are planned. \lname{Hyp.RamNrdSurj} (now discharged by
\lname{Hyp.ramNrdSurj\_holds}) enters every \lname{\_of\_SA\_HS} result at the standard levels through \lname{Tree.isQAdmissible\_*}. \lname{Hyp.InvLimExact} reaches
\lname{lim\_exact\_ribet} only through its conclusion (argument \lname{hA}, supplied as \lname{hI hconn});
\lname{Hyp.LambdaJL} is the only existing hypothesis reused in stage~6.}
\label{fig:hyp-deps}
\end{figure}

\begin{figure}[!tbp]\centering
\resizebox{\textwidth}{!}{% SPDX-License-Identifier: CC-BY-4.0
% Copyright (C) 2026 Fabian Januszewski, Abhijit (the Langlandscape team)
% Stage map / timeline: one row per hypothesis, a bar from the introducing stage, events marked.
% Colour = kind (Okabe--Ito); bar pattern repeats the kind for greyscale; event glyphs are shapes, not colours.
\begin{tikzpicture}[x=1.75cm,y=0.42cm,font=\sffamily\scriptsize]
  \tikzset{
    bar/.style={line width=4pt,line cap=butt},
    lit/.style={bar,draw=oiBlue!55},
    inf/.style={bar,draw=oiOrange!70,dash pattern=on 5pt off 1.5pt},
    oos/.style={bar,draw=oiVerm!65},
    con/.style={bar,draw=oiPurple!65,dash pattern=on 1.5pt off 1.5pt},
    dis/.style={bar,draw=oiGreen!70},
    pln/.style={line width=1pt,draw=oiSky!80!black,dashed},
    intro/.style={circle,fill=black,inner sep=1.4pt},
    corr/.style={regular polygon,regular polygon sides=3,fill=white,draw=black,inner sep=1.1pt,line width=0.6pt},
    disc/.style={rectangle,fill=oiGreen!80!black,draw=black,inner sep=2.1pt},
    plan/.style={circle,draw=oiSky!80!black,fill=white,inner sep=1.4pt,line width=0.8pt},
  }
  % stage columns: S0k at x = k
  \foreach \k in {1,...,6}{
    \draw[black!15] (\k-0.5,0.8) -- (\k-0.5,-20.6);
    \node[font=\sffamily\bfseries\scriptsize] at (\k,1.3) {S0\k};
  }
  \draw[black!15] (6.5,0.8) -- (6.5,-20.6);
  \node[font=\sffamily\tiny,text=black!60] at (4,1.95) {now (2026-11-15): stage 4 closing};
  \fill[black!6] (4.5,0.8) rectangle (6.5,-20.6);
  % rows: name / kind / start stage / end (6.45 = still open at end of project horizon)
  \def\row#1#2#3#4#5{% #1 y, #2 name, #3 style, #4 start, #5 end
    \node[anchor=east] at (0.45,#1) {\texttt{#2}};
    \draw[#3] (#4-0.35,#1) -- (#5,#1);
    \node[intro] at (#4-0.35,#1) {};}
  \row{0}{EichlerMass}{lit}{1}{6.45}
  \row{-1}{JLClassical}{oos}{1}{6.45}
  \row{-2}{LambdaJL}{oos}{1}{6.45}
  \row{-3}{SA}{lit}{2}{6.45}
  \row{-4}{HS}{lit}{2}{6.45}
  \row{-5}{BTtree}{lit}{2}{6.45}
  \row{-6}{SerreCoh}{lit}{2}{6.45}
  \row{-7}{SerreFree}{lit}{2}{6.45}
  \row{-8}{RamNrdSurj}{lit}{2}{6.15}
  \row{-9}{IharaH0}{inf}{3}{6.45}
  \row{-10}{EllipticCount}{lit}{3}{6.45}
  \row{-11}{BrandtStabFree}{inf}{3}{6.45}
  \row{-12}{PadicOrdDecomp}{dis}{4}{4.15}
  \row{-13}{MordLambdaFree}{inf}{4}{6.45}
  \row{-14}{HidaControl}{lit}{4}{6.45}
  \row{-15}{HsiehCR}{con}{4}{6.45}
  \row{-16}{InvLimExact}{inf}{4}{6.45}
  % planned rows
  \node[anchor=east] at (0.45,-17.5) {\texttt{X.char, X.gal, X.iso}};
  \draw[pln] (4.65,-17.5) -- (6.45,-17.5); \node[plan] at (4.65,-17.5) {};
  \node[anchor=east] at (0.45,-18.5) {\texttt{HsiehGor} (moved S05$\to$S06)};
  \draw[pln] (5.65,-18.5) -- (6.45,-18.5); \node[plan] at (5.65,-18.5) {};
  \node[anchor=east] at (0.45,-19.5) {\texttt{OldDecomp, \dots, KimOtaRT} (8)};
  \draw[pln] (5.65,-19.5) -- (6.45,-19.5); \node[plan] at (5.65,-19.5) {};
  % events: corrections (triangles) with short labels, discharge (square)
  \def\ev#1#2#3{\node[corr] at (#1,#2) {}; \node[font=\sffamily\tiny,anchor=south west,inner sep=1pt] at (#1+0.04,#2+0.05) {#3};}
  \ev{1.15}{-1}{hfin (L16c)}
  \node[corr] at (3.8,-2) {}; \node[font=\sffamily\tiny,anchor=south east,inner sep=1pt] at (3.76,-1.95) {$\langle d\rangle^{-1}$ (X1 / E1L)};
  \ev{4.4}{-2}{guard (L17a R1)}
  \ev{2.3}{-3}{finite index, Nikolov--Segal (L13a)}
  \ev{2.3}{-6}{citations (H2)}
  \ev{2.3}{-7}{citations (H2)}
  \ev{2.3}{-8}{vacuous instance removed (M1)}
  \node[disc] at (6.15,-8) {}; \node[font=\sffamily\tiny,anchor=south east,inner sep=1pt] at (6.1,-7.95) {discharged (final gate)};
  \ev{3.3}{-11}{status/route (L13)}
  \node[disc] at (4.15,-12) {}; \node[font=\sffamily\tiny,anchor=west] at (4.25,-12) {discharged: \texttt{padicOrdDecomp\_holds} (L3)};
  \ev{4.4}{-13}{scope note (L17c R14)}
  \ev{4.4}{-14}{reason text (L17a R4)}
  \ev{4.4}{-16}{connectivity premise (L17c R13)}
  \node[font=\sffamily\tiny,anchor=south west,inner sep=1pt] at (3.65,-8.95) {clause (b) deferred to S04, not done};
  % legend
  \begin{scope}[shift={(0.2,-21.8)}]
    \node[intro] at (0,0) {}; \node[anchor=west] at (0.08,0) {registered};
    \node[corr] at (1.2,0) {}; \node[anchor=west] at (1.28,0) {corrected};
    \node[disc] at (2.4,0) {}; \node[anchor=west] at (2.48,0) {discharged};
    \node[plan] at (3.6,0) {}; \node[anchor=west] at (3.68,0) {planned};
    \node[anchor=west] at (4.75,0) {bar = open; shaded = future stages};
  \end{scope}
\end{tikzpicture}
}
\caption{Stage map: when each hypothesis was registered, corrected and discharged (stages 1--4), and what is
planned (stages 5--6). Bar colour and pattern encode the kind, as in Figure~\ref{fig:hyp-deps}; labels name the work package or audit item; details in the entries.}
\label{fig:hyp-timeline}
\end{figure}

%======================================================================================================
\section{The hypotheses, one by one}\label{sec:hyp-entries}

Each entry gives: the informal statement; the Lean name and file (relative to \lname{lean/Langlandscape/}); the kind
and source; where it is used (Lean declarations, chapter labels); the current status; what would discharge it; and its
history. Chapter labels refer to \lname{mathematics/chapters/}.

%------------------------------------------------------------------------------------------------------
\subsection{Stage 1: targets and the mass formula}

\begin{hypentry}{GL2TF.Hyp.EichlerMass}{\kindlit}{S01 (reused from the import)}
\item[Statement] For a maximal order $\cO$ of $B=(a,b)_\Q$ that is ramified exactly at the primes dividing $D$:
  $\operatorname{mass}(\cO)=\sum_{[I]} 1/|\cO_\ell(I)^\times| = \frac{1}{24}\prod_{p\mid D}(p-1)$ (Eichler's mass formula,
  full unit groups). The ramification datum is an antecedent, so the family is true for every index.
\item[Lean] \lname{GL2TF.Hyp.EichlerMass a b D O} in
  \lname{trace-formula-project-import/lean/AlgJL/Hyp/EichlerMass.lean} (tag \lname{III.qcls.hyp:eichler-mass});
  alias \lname{Langlandscape.Hyp.EichlerMass} (\lname{Hyp/JLClassical.lean}, an \lname{abbrev}).
\item[Source] \emph{cited}: Eichler \cite{Eic55}; Shimizu \cite[(20)--(21), p.~180]{Shz65}; Pizer
  \cite[P0.5]{Piz80}; Voight \cite{Vo}. Chapter: \lname{thm:s01-mass}(a), \lname{if:s01-mass}.
\item[Used by] \lname{Quat.card\_XB\_Utower} (\lname{thm:s01-mass}(c),(d)), \lname{Smoke.eichlerMass\_apply},
  \lname{Ribet.sum\_inv\_wx\_Ohat}, \lname{Ribet.mazur\_check} (\lname{prop:s03-mazur}),
  \lname{Ordinary.vp\_card\_XB\_Utower\_ge} (stage 4). (Registry updated at S05-E1, item~\ref{inc:usedby}.)
\item[Status] Open in this project. Closed instances exist in the import's \emph{test} library
  (\lname{GL2TFTests.eichlerMass\_d11}, mass $5/12$; \lname{GL2TFTests.eichlerMass\_d2}, Hurwitz order, mass $1/24$),
  which is not a dependency of \lname{lean/Langlandscape/}; non-vacuity is therefore shown, but the instances cannot be
  used from our code.
\item[Discharge] In the import's plan: chapter III.13 (``Q8'') of the GL(2) trace-formula project, i.e.\ the
  Eichler mass formula from the trace formula / zeta functions of $B$. Not planned here.
\item[History] Unchanged since S01; reused, never redefined.
\end{hypentry}

\begin{hypentry}{Hyp.JLClassical}{\kindoos}{S01-L13}
\item[Statement] Classical Jacquet--Langlands for $k=2$, trivial coefficients: for $M$ prime to $D$ there is an
  isomorphism $\C^{X_B(U_0(M))}\cong S_2(\Gamma_0(DM))^{D\text{-new}}\oplus\C\cdot\mathrm{Eis}$ intertwining $T_\ell$ for
  all primes $\ell\nmid DM$ ($T_\ell$ acts on the Eisenstein line by $\ell+1$).
\item[Lean] \lname{Hyp.JLClassical d M hM S} in \lname{Hyp/JLClassical.lean}; the GL$_2$ side is abstract data
  \lname{S : GL2CuspData (D*M) 2} (a finite-dimensional $\C$-space with commuting operators), because Mathlib v4.35 has
  neither $S_k(\Gamma_0(N))$ with Hecke operators nor the comparison. With arbitrary \lname{S} the proposition only says
  that \lname{S} is a Hecke-module complement of the constants: it has no GL$_2$ content (\emph{schematic}, F6).
\item[Source] \emph{cited}: Jacquet--Langlands \cite{JL}; Gross \cite{Gr99}. Chapter: \lname{thm:s01-JLclassical}
  (\stcited; \lname{\textbackslash lean} mark ``schematic interface'').
\item[Used by] No Lean consumer; it is a target statement.
\item[Status] Open (schematic); no instance (none is feasible without GL$_2$ Hecke theory).
\item[Discharge] Not discharged during the event; planned for Stage~2 via the GL(2) trace-formula formalization (route~1 of
  \lname{ss:s00-whycited}), timing uncertain (see the report's Stage-2 plan). It would also remove the JL input to the comparison
  beyond $D=q$ (the Shimura-curve version of \lname{Hyp.X}), provided the trace-formula comparison is carried out for the quaternion algebra occurring there. It needs the classical JL correspondence formalized, e.g.\ via a complete GL(2) trace
  formula and its comparison with the quaternionic side (the import provides part of the trace formula, not the
  comparison), together with Hecke operators on $S_2(\Gamma_0(N))$ in Mathlib.
\item[History] S01 Lean audit R1/R9: the finiteness premise \lname{hfin} of $T_\ell$ was universally quantified
  and unprovable; discharged by \lname{relIndex\_elemT\_ne\_zero\_U0} (S01-L16c, \lname{6f229dc8}). Its docstring now cites the
  \lname{mathematics/} label (S05-E1, item~\ref{inc:planstaged}).
\end{hypentry}

\begin{hypentry}{Hyp.LambdaJL}{\kindoos}{S01-L13}
\item[Statement] $\Lambda$-adic Jacquet--Langlands: a surjection of $\Lambda$-algebras
  $\varphi:\T^{\mathrm{ord}}_{\mathrm{GL}_2}\to\T^{\mathrm{ord}}_B$ with kernel the $D$-new ideal, matching
  $T_\ell\mapsto T_\ell$ ($\ell\nmid DMp$), $U_\ell\mapsto U_\ell$ ($\ell\mid Mp$), $U_\ell\mapsto\Pi_\ell$ ($\ell\mid D$),
  and the diamonds by \emph{classical $\langle d\rangle\mapsto$ seed $\langle d\rangle^{-1}$}, for $d$ prime to $Mp$.
\item[Lean] \lname{Hyp.LambdaJL D M p Λ TGL TB genGL genB newIdeal} in \lname{Hyp/LambdaJL.lean}; both Hecke
  algebras, their generators and the new ideal are abstract (\emph{schematic}). The diamond clause is
  \lname{∀ d, Nat.Coprime d (M*p) → φ (genGL (.dia d)) * genB (.dia d) = 1}.
\item[Source] \emph{cited}: Chenevier \cite{Ch}; Longo--Vigni \cite[Prop.~6.1]{LV}; Hsieh \cite[\S4.5]{Hs}.
  Chapter: \lname{thm:s01-lambdaJL}, \lname{rem:s01-diamond}, \lname{lem:s01-Upconv}, \lname{if:s01-targets}.
\item[Used by] \lname{Defect.Setting.gor\_new}, \lname{Defect.defect\_formula}, \lname{Defect.cor\_allweights} (explicit
  argument \lname{hJL}; stage 6; only surjectivity and the kernel are used); otherwise a target statement, schematic.
\item[Status] Open (schematic); no instance.
\item[Discharge] Not during the event: needs classical JL plus Hida theory on the GL$_2$ side; its classical input could come
  from route~1 (\lname{ss:s00-whycited}; the Stage-2 plan: ``possibly the classical input of \lname{LambdaJL}''). Stage~06 takes it
  as a cited input because it does not use \lname{Hyp.X}.
\item[History, corrections] (1) \textbf{Inverse diamond (X1 erratum).} The original clause
  $\varphi(\langle d\rangle_{\mathrm{GL}})=\langle d\rangle_B$ encoded a false matching: the seed's adelic
  $\langle d\rangle$ is the classical $\langle d\rangle^{-1}$ (\lname{rem:s01-diamond}). Found by the blind cross-account
  review (\lname{mathematics/audits/2026-11-14\_s01-xacct.md}, item X1, \emph{major}; the S01 Lean audit had missed
  it); fixed in the chapter by S04-E1M (\lname{65b8f55c}) and in Lean by S04-E1L (\lname{6acc2ab9}), stated as
  $\varphi\langle d\rangle_{\mathrm{GL}}\cdot\langle d\rangle_B=1$. (2) \textbf{Coprimality guard.} The unguarded unit
  condition forced $0=1$ at the junk value $\langle 0\rangle=0$, making the hypothesis unsatisfiable (S04 Lean audit,
  R1); S04-L17a (\lname{7f984fa3}) added \lname{Nat.Coprime d (M*p)}. The chapter now states the guard (S05-E1, item~\ref{inc:guard}).
\end{hypentry}

%------------------------------------------------------------------------------------------------------
\subsection{Stage 2: strong approximation, norms, trees}

\begin{hypentry}{Hyp.SA}{\kindlit}{S02-L4a}
\item[Statement] Strong approximation for $B^1$ at $S=\{\infty,q\}$, in the shape (SA$_U$): for every $q$-admissible
  level $U$, $\hat B^1=B^1\,(B^1_q\times U^{1,(q)})$.
\item[Lean] \lname{Hyp.SA d q hq := ∀ U, Tree.IsQAdmissible d q hq U → Tree.SAat d q hq U} in \lname{Hyp/SA.lean}.
\item[Source] \emph{cited}: Voight \cite[Main Thm.~28.5.3, Cor.~28.5.12]{Vo}; Vign\'eras \cite[III.4.3]{Vigneras1980};
  for finite-index (not necessarily open) levels also Nikolov--Segal \cite[Thm.~1.1]{NS07}. Chapter: \lname{thm:s02-SA}.
\item[Used by] \lname{Tree.BGConnected\_of\_SA\_HS}, \lname{Tree.surjType0\_of\_SA\_HS},
  \lname{quotVEquiv\_of\_SA\_HS}, \lname{quotEEquiv\_of\_SA\_HS}, \lname{quotEquiv\_oO/\_tO\_of\_SA\_HS}
  (\lname{prop:s02-BG}, \lname{prop:s02-quot}); and \lname{Ribet.ribet\_exact\_of\_SA\_HS},
  \lname{ribet\_rank\_of\_SA\_HS}, \lname{ihara\_surj\_of\_SA\_HS}, \lname{ihara\_ker\_const\_of\_SA\_HS}
  (\lname{thm:s03-ribet}, \lname{cor:s03-ihara}); the last four are missing from the registry.
\item[Status] Open; no instance. Per level it is replaced by the connectivity premise \lname{BGConnected}
  (\lname{prop:s02-BG}), which the \lname{\_uncond} stage-4 results take directly.
\item[Discharge] Strong approximation for $\mathrm{SL}_1(B)$ over $\Q$ (plan item L02.SA, post-hackathon); needs adelic
  algebraic groups in Mathlib. The import has a \lname{strong\_approximation} for GL$_2$ that does not cover $B^1$.
\item[History] S02 Lean audit: \lname{IsQAdmissible} says ``finite index'' where the tex says ``compact open'';
  the openness step via Nikolov--Segal was added to the docstring (S02-L13a, \lname{faf256dd}). For the project's levels
  ($\hat\cO^\times$, $U_0(N)$, $U_1(N)$, $U_s$, all open) Nikolov--Segal is not needed.
\end{hypentry}

\begin{hypentry}{Hyp.HS}{\kindlit}{S02-L4a}
\item[Statement] Hasse--Schilling for the definite algebra: every positive rational is the reduced norm of a unit of
  $B$, $\Nrd(B^\times)=\Q_{>0}$ (the inclusion $\subseteq$ is proved).
\item[Lean] \lname{Hyp.HS d := ∀ r : ℚ, 0 < r → ∃ γ : (QuatAlg ℚ d.a d.b)ˣ, nrd γ = r} in \lname{Hyp/HS.lean}.
\item[Source] \emph{cited}: Voight \cite[Main Thm.~14.7.4]{Vo}; Vign\'eras \cite[III.4.1]{Vigneras1980}. Chapter:
  \lname{lem:s02-normone}(b).
\item[Used by] \lname{Tree.XBone\_toXB\_bijective} and all \lname{\_of\_SA\_HS} results listed under \lname{Hyp.SA}.
\item[Status] Open; no instance.
\item[Discharge] Hasse--Schilling, i.e.\ the Hasse--Minkowski theorem for the ternary norm form (post-hackathon);
  needs Hasse--Minkowski over $\Q$ in Mathlib.
\item[History] Unchanged.
\end{hypentry}

\begin{hypentry}{Hyp.BTtree}{\kindlit}{S02-L5b}
\item[Statement] The coset graph on $\GL_2(\Q_q)/\Q_q^\times\GL_2(\Z_q)$ is a tree in which every vertex has $q+1$
  neighbours (the Bruhat--Tits tree). Only connectedness is used downstream.
\item[Lean] \lname{Hyp.BTtree q := (Tree.graph q).IsTree ∧ ∀ v, Nat.card (neighborSet v) = q + 1} in
  \lname{Hyp/BTtree.lean}.
\item[Source] \emph{cited}: Serre \cite[II.1.1, Thm.~1]{Se}; formalized (Lean~4, older Mathlib) by Ludwig--Merten
  \cite{LM}. Chapter: \lname{lem:s02-tree}(a).
\item[Used by] \lname{Tree.BGConnected\_of\_surj}, \lname{Tree.BGConnected\_of\_SA\_HS} (step (i)$\Rightarrow$(iii) of
  \lname{prop:s02-BG}) and the four \lname{Ribet.*\_of\_SA\_HS} (all take \lname{hBT}); the registry now lists them.
\item[Status] Open; no instance.
\item[Discharge] Port of the Ludwig--Merten formalization from Lean v4.19 to v4.35 (post-hackathon).
\item[History] S02 audit H6: docstring records that $(q+1)$-regularity is not consumed.
\end{hypentry}

\begin{hypentry}{Hyp.SerreCoh}{\kindlit}{S02-L3}
\item[Statement] If $\BG(U,q)$ is connected and all stabiliser orders $|\Gamma_x|$ are units in $R$, then
  $H_1(\Gamma_U,R)\cong\ker\partial$ and $H_0(\Gamma_U,R)\cong\coker\partial$ for the two-term complex $\partial$ of the
  quotient graph (trivial coefficients).
\item[Lean] \lname{Hyp.SerreCoh d q hq U hU hconn R} in \lname{Hyp/SerreCoh.lean} (Mathlib \lname{groupHomology};
  $R$ in \lname{Type} for universe reasons; the isomorphisms are abstract, not asserted compatible with $\partial$).
\item[Source] \emph{cited}: Serre \cite[II.2.8, Prop.~13 and Rem.~1]{Se}; Brown \cite[VII.7]{Br}. Chapter:
  \lname{prop:s02-coh}(b) (\stpmc; the tex has coefficients $V$, the Lean form trivial coefficients).
\item[Used by] Statement only (test \lname{LanglandscapeTests/S02/SerreShape.lean}).
\item[Status] Open; no instance.
\item[Discharge] Bass--Serre theory (homology of groups acting on trees) in Mathlib (post-hackathon, plan item L02.6).
\item[History] S02 audit H2/H6: citations and docstring corrected.
\end{hypentry}

\begin{hypentry}{Hyp.SerreFree}{\kindlit}{S02-L3}
\item[Statement] If $\BG(U,q)$ is connected and all $\Gamma_x(U)$ are trivial, then $\Gamma_U$ is free of rank
  $(q-1)|X(U)|+1$.
\item[Lean] \lname{Hyp.SerreFree d q hq U hU hconn hfree} in \lname{Hyp/SerreFree.lean} (via
  \lname{FreeGroupBasis}).
\item[Source] \emph{cited}: Serre \cite[I.3.3, Thm.~4$'$; II.1.5, Thm.~4]{Se}; Bass \cite[3.6]{Bass1993}.
  Chapter: \lname{prop:s02-free}(a).
\item[Used by] Statement only (test \lname{SerreShape.lean}).
\item[Status] Open; no instance.
\item[Discharge] Needs \lname{Hyp.BTtree} and Serre's theorem that a group acting freely on a tree is free
  (Bass--Serre theory), post-hackathon.
\item[History] S02 audit H2/H6: citations corrected.
\end{hypentry}

\begin{hypentry}{Hyp.RamNrdSurj}{\kinddis}{S02-L10}
\item[Statement] At every ramified prime $\ell\mid D$, the reduced norm maps the units of the local maximal order onto
  $\Z_\ell^\times$: $\Nrd(\cO_\ell^\times)=\Z_\ell^\times$.
\item[Lean] \lname{Hyp.RamNrdSurj d} in \lname{Hyp/RamNrdSurj.lean}.
\item[Source] \emph{cited}: Voight \cite[Lemma~13.4.9]{Vo}; Vign\'eras \cite[II.1.5]{Vigneras1980}; also proved on paper in
  \lname{lem:s01-nrd} from the cited local structure (Voight Thm.~13.3.11); \emph{formalized} (below). Chapter:
  \lname{lem:s01-nrd} (the $\ell\mid D$ ingredient; status reconciled across s01, s02 and Lean at S05-E1, item~\ref{inc:nrd}).
\item[Used by] \lname{Quat.exists\_Ohat\_nrd\_eq}, \lname{Quat.nrdSurj\_Ohat}, \lname{nrdSurj\_U0}, \lname{nrdSurj\_U1},
  \lname{nrdSurj\_Utower}; \lname{Tree.isQAdmissible\_Ohat}, \lname{\_U0}, \lname{\_U1}, \lname{\_Utower}. Hence, at the
  standard levels, every \lname{\_of\_SA\_HS} consumer.
\item[Status] \textbf{Discharged} by \lname{Langlandscape.Hyp.ramNrdSurj\_holds (d : DefiniteData) : Hyp.RamNrdSurj d}
  (\lname{Hyp/RamNrdSurjDischarge.lean}), via the reduction \lname{Hyp.ramNrdSurj\_of\_local} and the local inputs
  \lname{Local.exists\_nrd\_eq\_unit\_of\_ne\_two} (\lname{Local/QuatNrdOdd.lean}) and \lname{Local.exists\_nrd\_eq\_unit\_two}
  (\lname{Local/QuatNrdTwo.lean}); all four use only \lname{propext}, \lname{Classical.choice}, \lname{Quot.sound}
  (final gate, TITAN \lname{ramnrd-final-2} at \lname{2323a252}). The consumers keep the explicit argument and receive it
  as \lname{ramNrdSurj\_holds d}. Needed for every definite algebra ($D>1$ always).
\item[Discharge] Done (2026-11-15) by a route different from the chapter's: the local inputs show that every
  $w\in\Z_\ell^\times$ is a reduced norm from $(a,b\mid\Q_\ell)$ at \emph{every} prime $\ell$ and for all nonzero $a,b$
  (\emph{formalized}, the two \lname{Local} theorems); $\ell\mid D$ enters only through $\cO_\ell=\{x:\Nrd x\in\Z_\ell\}$.
  The proof was drafted by an external model (mistral-large-4-0 via opencode), outside our agent process; see the
  final report, \S3.
\item[History] Introduced as a deviation of S02-L10. Its only instance, for $D=1$, was vacuous (no definite algebra
  has $D=1$; S02 audit M1) and was deleted (S02-L13a, \lname{faf256dd}). Discharged after the S06 gate (merge
  \lname{2323a252}; independent Fable audit PASS WITH REMARKS).
\end{hypentry}

%------------------------------------------------------------------------------------------------------
\subsection{Stage 3: Ribet's sequence, Ihara, Mazur}

\begin{hypentry}{Hyp.IharaH0}{\kindinfra}{S03-L11}
\item[Statement] Ihara's lemma with coefficients, clause (a): under the standing hypotheses of \lname{ss:s03-ihara}
  ($p\ge5$, $p\ne q$, $p\nmid D$, $U$ $q$-admissible and factorisable at $p$, $\BG(U,q)$ connected, $V$ a f.g.\
  $\Z_p$-module with continuous $U_p$-action, $-1$ acting trivially), the map
  $\varepsilon:\Brandt{U}{V}^2\to V_{U^1_p}$ is well defined and onto, and $\ker\varepsilon=\operatorname{im}\alpha$ for
  $\alpha=(-s_{0*},s_{1*})$; hence $\coker\alpha\cong V_{U^1_p}$.
\item[Lean] \lname{Hyp.IharaH0 d q hq U hU p hp V c} in \lname{Hyp/IharaH0.lean}.
\item[Source] \emph{proved modulo cited} in the chapter (\lname{thm:s03-iharaH0}(a), \stpmc): well-definedness,
  $\varepsilon\circ\alpha=0$ and surjectivity are elementary; only $\ker\varepsilon\subseteq\operatorname{im}\alpha$
  uses the density of $\Gamma_U$ in $U^1_p$ (full strong approximation, \lname{cor:s03-ihara}(ii$'$)). Classical
  analogue: Diamond--Taylor \cite[pp.~445--446]{DT}.
\item[Used by] Statement only (tests; intended for stage-6 mathematics and the talk).
\item[Status] Open; no instance. Premises jointly satisfiable at $\hat\cO^\times$ with trivial $\rho$ (S03 Fable
  audit: not vacuous).
\item[Discharge] Prove the elementary conjuncts now and reduce to $\ker\varepsilon\le\operatorname{range}\alpha$; that
  part needs \lname{Hyp.SA} and the cohomological interface \lname{if:s02-coh}. Post-hackathon.
\item[History] Clause (b) (Hecke action) excluded from the start and deferred as an optional output of S04-L15a,
  which merged without it; the Hecke operators with coefficients ``from S04-L15a'' now exist
  (\lname{Brandt/HeckeCoeff.lean}, \lname{Brandt/UpCoeff.lean}); registry updated at S05-E1 (item~\ref{inc:ihara}), clause (b) still not stated.
\end{hypentry}

\begin{hypentry}{Hyp.EllipticCount}{\kindlit}{S03-L6}
\item[Statement] For $D=q$ prime and $U=\hat\cO^\times$: if $q\ge5$, every class has $w_x\in\{1,2,3\}$, exactly
  $[q\equiv3\bmod4]$ classes have $w_x=2$ and exactly $[q\equiv2\bmod3]$ classes have $w_x=3$; if $q\in\{2,3\}$, then
  $|X_B(\hat\cO^\times)|=1$.
\item[Lean] \lname{Hyp.EllipticCount d} in \lname{Hyp/EllipticCount.lean}.
\item[Source] \emph{cited}: Voight \cite[Thm.~11.5.14, 25.1.1, 30.1.5]{Vo}; Mazur \cite[App., Thm.~1.1(d)]{Mz}.
  Chapter: \lname{lem:s03-elliptic} (\stpmc) and the $q\in\{2,3\}$ part of \lname{prop:s03-mazur}.
\item[Used by] \lname{Ribet.mazur\_check} (with \lname{EichlerMass}; every prime $D$, the $q\le3$ clause included, so the $q=2,3$ Mazur check is formalized modulo this hypothesis; item~\ref{inc:mazur}); tests \lname{test\_mazur\_check\_*}.
\item[Status] Open; no closed instance. True for every prime $D$, hence non-vacuous.
\item[Discharge] The Eichler class number formula and the classification of unit groups of maximal orders; both are
  absent from Mathlib and the import.
\item[History] Cross-account review s03 X6 (Lean stronger than the lemma, because of the $q\in\{2,3\}$ clause):
  resolved by the chapter note.
\end{hypentry}

\begin{hypentry}{Hyp.BrandtStabFree}{\kinddis}{S03-L8}
\item[Statement] At $U=\hat\cO^\times$, for a prime $\ell\nmid D$ and classes $c\ne c'$, the weight $w_c$ divides the
  Brandt matrix entry $B(\ell)_{cc'}$ (the stabiliser $\Gamma_c/\{\pm1\}$ acts freely on the neighbours in a different
  class).
\item[Lean] \lname{Hyp.BrandtStabFree d ℓ} in \lname{Hyp/BrandtStabFree.lean}; it also assumes $D$ prime, which
  the proof does not use, and is stated only at $\hat\cO^\times$ (the tex holds for all levels $U_0(M)$).
\item[Source] \emph{proved} in the chapter (\lname{lem:s03-stabfree}, \stproved, elementary: $\Z[\gamma]$ is $\Z[i]$ or
  $\Z[\omega]$, Euclidean, so $\pi\in\Z[\gamma]$ with $\Nrd\pi=\ell$). Cross-reference: Ribet \cite[Thm.~3.12]{Ri90}.
\item[Used by] \lname{Ribet.phiComb\_eisenstein\_Ohat}, \lname{Ribet.phiComb\_eisenstein\_DR},
  \lname{Ribet.brandtZ\_dvd} (step (1) of \lname{prop:s03-phieis}).
\item[Status] \textbf{Discharged} by \lname{Langlandscape.Hyp.brandtStabFree\_holds}
  (\lname{Hyp/BrandtStabFreeDischarge.lean}): an external draft (mistral-large-4-0 via opencode, on the second team member's machine,
  outside our agent process), integrated into the library, independently audited (Fable, PASS) and gated by our agents.
\item[Discharge] (Plan before the discharge.) The neighbour/lattice description (\lname{lem:s03-nbr}) of the import's Brandt matrix plus the
  arithmetic of $\Z[i],\Z[\omega]$ inside $B$. Working estimate (item~\ref{inc:brandt}): 1000--1500 lines via Morita theory plus a PID argument (probed
  2026-11-14); the 500--800 line elementary route in \lname{STATE.md} is unprobed. The unused premise ``$D$ prime'' is
  recorded in the registry (open design decision). Deferred post-stage-04.
\item[History] Registered with a ``Morita + PID'' source; the S03 Lean audit R1 found the chapter proof elementary;
  status and route corrected in S03-L13 (\lname{e63c5c48}, \lname{fae211a3}).
\end{hypentry}

%------------------------------------------------------------------------------------------------------
\subsection{Stage 4: Hida theory, control, the $\Lambda$-adic Ribet sequence}

\begin{hypentry}{Hyp.PadicOrdDecomp}{\kinddis}{S04-L1}
\item[Statement] Every endomorphism $T$ of a finitely generated $\Z_p$-module $M$ has an ordinary decomposition
  $M=M^{\mathrm{ord}}\oplus M^{\mathrm{nil}}$ ($T$ invertible on the first summand, topologically nilpotent on the second).
\item[Lean] \lname{Hyp.PadicOrdDecomp p} in \lname{Hyp/PadicOrdDecomp.lean}.
\item[Source] \emph{proved} in the chapter (\lname{def:s04-e}, Step~5) and \emph{formalized}.
\item[Used by] About 25 tower theorems in \lname{Ordinary/\{RibetFinite, RibetFiniteFull, Tower, TowerRank, Limit\}.lean}
  as the argument \lname{hdec}, including \lname{Ordinary.lim\_exact\_ribet}, the L10b theorem \lname{ordFree\_finite} and the L17b
  \lname{RibetFiniteFull} results (item~\ref{inc:usedby}).
\item[Status] \textbf{Discharged} by \lname{Langlandscape.Ordinary.padicOrdDecomp\_holds}
  (\lname{Ordinary/PadicDecomp.lean}, S04-L3, Weierstrass-preparation route; standard axioms only, gate
  \lname{s04-int4b-1}). At \lname{cd0a268} every consumer has an \lname{\_uncond} twin in
  \lname{Ordinary/TowerUncond.lean} (18 theorems, e.g.\ \lname{lim\_exact\_ribet\_uncond},
  \lname{ordRibet\_exact\_Utower\_uncond}, \lname{ord\_control\_finite\_uncond}).
\item[History] Registered at S04-X1 as ``to be proved'' (fallback F-e of the stage plan); discharged in S04-L3
  (\lname{a2326c23}); stale ``to be discharged'' text fixed in L17a (R4); \lname{\_uncond} forms completed in L17b/L17c
  (R9, R11).
\end{hypentry}

\begin{hypentry}{Hyp.InvLimExact}{\kindinfra}{S04-L14b}
\item[Statement] If $\BG(UU_s,q)$ is connected for all $s\ge1$, the transition maps
  $e\,H_1(\Gamma_{UU_{s+1}})\to e\,H_1(\Gamma_{UU_s})$ of the $A$-system of the ordinary Ribet sequence are surjective.
\item[Lean] At \lname{cd0a268}: \lname{Hyp.InvLimExact H := (∀ s ≥ 1, Tree.BGConnected \dots (d.Utower M p s \dots))
  → ∀ n, Function.Surjective (H.transA n)} in \lname{Hyp/InvLimExact.lean}. At \lname{d214fb58} the connectivity
  premise is absent.
\item[Source] \emph{proved} in the chapter (\lname{thm:s04-lambdaribet}(a), \stproved) by the D-lim route: control at
  both levels, $\Delta$-linearity of $\alpha$, snake lemma, under connectivity. The compactness proof in the same theorem
  gives exactness of the limit \emph{without} this surjectivity, so it does not prove this statement.
\item[Used by] \lname{Ordinary.lim\_exact\_ribet} and \lname{lim\_exact\_ribet\_uncond}, but only through their
  argument \lname{hA : ∀ n, Function.Surjective (H.transA n)}, which is supplied as \lname{hI hconn} from
  \lname{hI : Hyp.InvLimExact H} (test \lname{LanglandscapeTests/S04/HypLam.lean}). No library declaration mentions
  \lname{Hyp.InvLimExact} in its type (registry row and chapter corrected at S05-E1, item~\ref{inc:invlim}).
\item[Status] Open (the conditional fallback of S04-L13b); no instance (needs a concrete \lname{RibetHyp}); not
  vacuous.
\item[Discharge] Formalize the D-lim snake route on the concrete Ribet systems: the abstract snake step
  \lname{Ordinary.surjective\_of\_snake} exists; missing are the control kernel $\ker=I\cdot B$ and the
  $\Delta$-equivariance of $\alpha$ (deviation 7 dropped this in stage 4). Connectivity itself follows from
  \lname{Hyp.SA}, \lname{Hyp.HS}, \lname{Hyp.BTtree} (cross-account review s04, X8).
\item[History, corrections] Introduced after the stage-4 plan audit (Muse, F2) found the fallback undeclared.
  \textbf{Connectivity premise added} by S04-L17c (\lname{3013fb81}, audit R13): the tex assumes $\BG(UU_s,q)$ connected
  (global seam audit F7, decision D5; \lname{s04/lambdaribet.tex}, marked \lname{\textbackslash provisional}), and the
  old ``compactness'' justification was wrong. Registry row and chapter updated at S05-E1 (item~\ref{inc:invlim}).
\end{hypentry}

\begin{hypentry}{Hyp.MordLambdaFree}{\kindinfra}{S04-L14b}
\item[Statement] $\Mord(U)=\varprojlim_s e\,\Z_p[X_s(U)]$ is a free $\Lambda$-module of finite rank, and for every $n$
  the projection to level $n+1$ is onto with kernel $\ker(\Lambda\to\Z_p[\Delta_{n+1}])\cdot\Mord(U)$
  (control, $\Lambda$-adic half).
\item[Lean] \lname{Hyp.MordLambdaFree T π ρ hρ} in \lname{Hyp/MordLambdaFree.lean}; the $\Lambda$-module structure
  comes from explicit levelwise diamond data \lname{(ρ, hρ)}. It is the theorem \emph{only} for the intended data
  ($G_n=\Delta_{n+1}$, reduction maps, diamond action); these data do not yet exist in Lean (scope note, R14,
  \lname{cd0a268}). The rank is $\exists r$; Hecke compatibility is not part of the statement.
\item[Source] Chapter \lname{thm:s04-control}(a): \stproved{} (raised from \stpmc{} after S04-M9); also
  \emph{cited}: Hsieh \cite[Thm.~4.2]{Hs}; Longo--Vigni \cite[Prop.~3.6]{LV11}. Lean docstring and registry now say ``tex proved'' (S05-E1, item~\ref{inc:mord}).
\item[Used by] Statement only (test \lname{HypLam.lean}); no consumer is named in the stage-5/6 plans (grep at S05-E1).
\item[Status] Open; no instance (a vacuous one is forbidden); toy case $\Delta_s=1$ described in the docstring only.
\item[Discharge] Build the diamond data, compatible bases $e\,\Z_p[X_s]\cong\Z_p[\Delta_s]^r$ (Nakayama), and the
  Iwasawa algebra; post-stage.
\item[History] Scope note added by L17c (R14).
\end{hypentry}

\begin{hypentry}{Hyp.HidaControl}{\kindlit}{S04-L14a}
\item[Statement] Weight control, definite side, finite level mod $p^s$: for $k\ge2$ and $U=U_1(M)$-type levels with
  $p\nmid\varphi(M)$, the ordinary part of the weight-$k$ Brandt module mod $p^s$ is isomorphic to the ordinary part of
  $e\,\cO[X_s(U)]/P_k$ (\lname{def:s04-Pk}).
\item[Lean] \lname{Hyp.HidaControl M p s k hk hφ A Mk Uk M2 U2 Pk Qk π Ubar} in \lname{Hyp/HidaControl.lean};
  \emph{schematic}: the Hecke modules are abstract data (decision Q3, taken at int3 before the coefficient Hecke
  operators \lname{Brandt/HeckeCoeff.lean}, \lname{UpCoeff.lean} were green). Stated on modules, not on their duals as
  in the tex.
\item[Source] \emph{cited}: Hsieh \cite[Thm.~4.2(2)]{Hs}; Longo--Vigni \cite[Thm.~3.7]{LV11}; Hida
  \cite[Thm.~9.4]{Hi88b}. Chapter: \lname{thm:s04-control}(b), \lname{rem:s04-controlb}.
\item[Used by] Statement only (test).
\item[Status] Open (schematic); no instance (an instance on abstract data would be vacuous).
\item[Discharge] A faithful statement first: $U_p$ on $V_k$-coefficients at all levels $U_s$ (\lname{UpV}, S04-L15b),
  the quotient by $P^{\mathrm{ad}}_k$, and $\Mord$. Cited, post-hackathon.
\item[History] S04 audit R4: the stated reason (``coefficients not available'') was stale; corrected in L17a.
  Open: audit item O1 (\lname{hk} unused; now recorded in registry and chapter, not applied).
\end{hypentry}

\begin{hypentry}{Hyp.HsiehCR}{\kindcond}{S04-L14a}
\item[Statement] Hsieh's condition (CR$_D$) on a residual representation $\bar\rho$: absolutely irreducible,
  $p$-distinguished ($\chi^+\ne\chi^-$), and ramified at every $\ell\mid D$ with $\ell\equiv1\bmod p$.
\item[Lean] \lname{Hyp.HsiehCR p D G absIrred chiPlus chiMinus ram} in \lname{Hyp/HsiehCR.lean}; the residual
  representation is abstract (\lname{absIrred}, \lname{ram} are data).
\item[Source] \emph{cited}: Hsieh \cite[\S1.4, Hyp.~(CR,$\Sigma$)]{Hs}. Chapter: \lname{hyp:s04-CR} (\lname{\textbackslash lean} mark added at S05-E1, item~\ref{inc:cr}).
\item[Used by] None. The stage-6 plan's \lname{Hyp.HsiehGor} has no CR premise and does not mention \lname{HsiehCR}; the registry row now says so.
\item[Status] Open; a condition, never discharged. Toy instance \lname{LanglandscapeTests.S04.Hyp.hsiehCR\_toy}
  ($p=5$, $D=3$; audit O2: clause (iii) is not exercised).
\item[Discharge] None; to be checked for each $\bar\rho$ (e.g.\ in the HACS examples).
\end{hypentry}

%======================================================================================================
\section{The imported trace-formula fork}\label{sec:hyp-import}

\lname{lean/trace-formula-project-import/} is a curated fork of the GL(2)/$\Q$ trace-formula and quaternion-algebra
material of an earlier project, with Lean packages \lname{GL2TF}/\lname{AlgJL} (library roots \lname{AlgJL},
\lname{AlgJLTests}), LaTeX, and a LaTeX back-translation. It is frozen (\lname{agents/rules/frozen.md}).
\lname{lean/Langlandscape/} depends on it via its \lname{lakefile.toml} (\lname{require GL2TF}); our code imports
35 \lname{AlgJL.*} modules (transitive closure, \emph{computed} by grep).

\paragraph{Hypotheses of the import.} Exactly one, \lname{GL2TF.Hyp.EichlerMass} (first entry of \S\ref{sec:hyp-entries}
above; import README; \lname{latex/shared/hypotheses.tex}). The import follows the same policy: no axioms, hypotheses
as \lname{def \dots : Prop} arguments, the standard axioms only.

\begingroup\sloppy
\paragraph{The \lname{sorry}s of the import.} The brief for this document mentions one; there are \emph{two}
(import README, ``Open sorrys: two''), both in the main library \lname{AlgJL}, both stretch examples already open in
the parent project:
\begin{enumerate}[leftmargin=1.5em]
\item \lname{GL2TF.Quat.eichler\_mass\_instances\_small}, \lname{AlgJL/VolIII/P2/Instances.lean:1815} (declaration line;
  the \lname{sorry} token is at line 1832), tag \lname{III.qcls.ex:eichler-mass-instances-small}: for Pizer's orders with
  $(a,b,D)=(-1,-3,3),(-2,-5,5),(-1,-7,7),(-2,-13,13)$ the order is maximal, has one class, unit group of order
  $12,6,4,2$, and satisfies \lname{Hyp.EichlerMass}.
\item \lname{GL2TF.Quat.hecke\_d11\_weight4}, \lname{AlgJL/VolIII/Brandt/BrandtInstances.lean:783} (declaration; token
  at 794): $\dim M_4(\cO_{11})=2$ with $\operatorname{tr}T_2=2$, $\det T_2=-2$.
\end{enumerate}
\emph{What depends on them:} no Lean declaration refers to either (they are leaves; only the B-edition LaTeX cites them).
The registry instances \lname{eichlerMass\_d11}, \lname{eichlerMass\_d2} are derived from other, proved theorems.
\emph{Do our theorems use them?} No. Neither module is in the import closure of the library \lname{Langlandscape}.
\lname{AlgJL.VolIII.P2.Instances} is imported by two test files,
\lname{LanglandscapeTests/S01/CentralSimple.lean} and \lname{LocalComponent.lean}, which do not use the sorried theorem;
this is why the gate build log shows one ``declaration uses sorry'' warning, at \lname{Instances.lean:1815}, and none for
\lname{BrandtInstances.lean}, which is never built. \emph{What the clean audit means precisely:} for each of the 592
audited declarations, \lname{sorryAx} is not among the axioms of its transitive dependency cone; so none of them depends,
however indirectly, on either sorried theorem. It does \emph{not} say that the import is sorry-free, and it says nothing
about declarations outside the audit list (e.g.\ tests).

\par\endgroup
\paragraph{The import's own audit.} \lname{latex/backtranslation/generated/axioms.tex} (frozen from parent commit
\lname{84818d0}, not re-run in the fork) has 525 entries; exactly the two above contain \lname{sorryAx}.

\paragraph{Vendored Apache-2.0 files.} Vendored files keep their upstream headers (\lname{agents/rules/licensing.md}).
\begin{itemize}[leftmargin=1.2em]
\item In the import, \lname{lean/Vendor/} (\lname{MANIFEST.tsv}, 4 files): from the Anthropic FLT repository
  \cite{FLTAnthropic} (\lname{6e837e753555}) \lname{Def\_LocalLanglands\_HeckePair.lean} (Hecke pairs, used by
  \lname{AlgJL/Common/HeckeModule.lean}) and \lname{Def\_LocalLanglands\_GelfandInvolution.lean} (used by
  \lname{AlgJL/VolI/P1/Spherical.lean}), imports adjusted only; from Imperial's FLT project \cite{FLTImperial}
  (\lname{1ab58c167e4f}) \lname{FLT/NumberField/Padics/RestrictedProduct.lean} and
  \lname{FLT/Mathlib/Topology/Algebra/RestrictedProduct/Basic.lean} (restricted products, the adeles of
  \lname{AlgJL/VolI/P0/Adeles.lean}), verbatim.
\item In our project, \lname{lean/Langlandscape/Vendor/AxiomStacks/} (8 files, verbatim, blob SHAs checked):
  Stacks-project material formalized by AxiomMath \cite{AxiomStacks} (\lname{c496df0d}): Skolem--Noether
  (\lname{StacksCore.RingTheory.SkolemNoether}, Stacks tags 074Q/074R) and its supporting files on simple rings, Morita
  equivalence for simple rings, two-sided ideals, module categories and finite dimension. Our only direct importer is
  \lname{Langlandscape/Quat/CentralSimple.lean} (\lname{QuatAlg.algHom\_conj}, \lname{QuatAlg.algEquiv\_inner},
  \lname{def:s01-B}).
\end{itemize}
Vendored files are proofs, not hypotheses; they are covered by the same axiom audit.

%======================================================================================================
\section{Hypotheses of stages 5 and 6}\label{sec:hyp-planned}

From \lname{plan/implementation/stage-05.md} (\S0, registry table), \lname{stage-06.md} (\S0) and
\lname{REPLAN-S02-06.md}. Both sets are now registered in Lean (\lname{HYPOTHESES.md} at \lname{99f5500f}); the stage-6 rows are also in
Table~\ref{tab:hyp-status}. The tables below record the plan and the current classification. No work package discharges an existing
hypothesis.

\paragraph{Stage 5 (the character group $\mathcal X$ and its $\Lambda$-adic version).} Proofs of \lname{Hyp.X} are
out of scope (stage-05 plan, \S0).
\begin{center}\small
\begin{tabular}{@{}>{\raggedright\arraybackslash}p{2.4cm}>{\raggedright\arraybackslash}p{5.2cm}>{\raggedright\arraybackslash}p{3.4cm}>{\raggedright\arraybackslash}p{3.4cm}@{}}\toprule
planned name & statement & source & consumer / instance\\\midrule
\lname{Hyp.X.char} & \lname{hyp:s05-X} (X1)--(X3): the character groups of the Jacobians at $q$, as one compatible
family $X_s\cong\deg^0(\Z[\text{supersingular}])$ & Ribet \cite[Prop.~3.1, 3.7, 3.8]{Ri90}; Deligne--Rapoport
\cite[V.1.16, 1.18]{DR}; Raynaud & S05-L7 \lname{xlambda\_finite}; toy \lname{canonicalCharData\_char}\\
\lname{Hyp.X.gal} & (X4): graded pieces of the Tate module, Frobenius clause in the $\mu_N$-model (FD1 decided, Fabian, 2026-11-15) & SGA7 IX
\cite{SGA7}; Illusie \cite[Thm.~2.7]{Ill}; Ribet--Stein \cite[L.3.17]{RS} & toy \lname{toyTate\_gal}\\
\lname{Hyp.X.iso} & \lname{def:s05-Xiso}, on abstract Hecke-module data & Hsieh \cite[Thm.~4.5]{Hs}; Helm
\cite[Cor.~8.11]{Hel} & optional S05-L8; toy \lname{toyXIso\_holds}\\\bottomrule
\end{tabular}
\end{center}
For $D=q$ the content of \lname{Hyp.X} is expected by the Deligne--Rapoport and Deuring geometry of the $\Gamma_0$ case, without JL
(\lname{rem:s05-Xstatus}, expected; not carried out). Beyond $D=q$, JL enters at the Shimura-curve comparison; a JL-free (geometric)
discharge there is open (route~2 of \lname{ss:s00-whycited}); route~1 removes that JL input through \lname{Hyp.JLClassical},
provided the trace-formula comparison is carried out for the quaternion algebra occurring there.
Reused, not redefined: \lname{GL2TF.Hyp.EichlerMass}, the stage-4 Hsieh/$\Mord$ hypotheses; S05-L7 takes
\lname{hdec : Hyp.PadicOrdDecomp}, which is discharged. Decision FD1 (OVERVIEW \S9, Q-A) was taken on 2026-11-15 ($\mu_N$-model with Picard functoriality).

\paragraph{Stage 6 (Gorenstein property and the defect formula).} Nine new files (S06-L8a, registered), each a
\lname{def \dots : Prop} over \lname{Defect.Setting}:
\begin{center}\small
\begin{tabular}{@{}>{\raggedright\arraybackslash}p{3.1cm}>{\raggedright\arraybackslash}p{5.2cm}>{\raggedright\arraybackslash}p{3.3cm}>{\raggedright\arraybackslash}p{2.6cm}@{}}\toprule
name & statement & source & discharge\\\midrule
\lname{Hyp.HsiehGor} & (Gor$_R'$), definite side: $\T_B$ Gorenstein and $M\cong\T_B$ (\lname{thm:s04-Gor}, stated over $\mathbb I$;
\lname{Setting.gor\_new} uses the $\Lambda$-form, \lname{Setting.gor\_new\_of\_I} the $\mathbb I$-form) & Hsieh
\cite[Thm.~4.5]{Hs} & cited\\
\lname{Hyp.OldDecomp} & \lname{lem:s06-decomp} (corrected new/old decomposition) & Hida \cite[Thm.~3.1, Cor.~3.2]{Hi86b} &
maths only\\
\lname{Hyp.HsiehGorFull} & (Gor$_R$) on the full Hecke algebra & \cite[\S3.3]{Hs}; \cite{Wi95}; \cite{MW86};
\cite[Cor.~8.11]{Hel} & cited\\
\lname{Hyp.FreeRankTwo} & inputs (1)--(3) and steps 1--3 of \lname{prop:s06-galois}: a Galois-theoretic input in hypothesis form;
given \lname{Hyp.ResidualBurnside} equivalent to $\tq_q\subseteq\cc_q$ (both directions in Lean) & Hsieh; tex & cited; on paper \stpmc\\
\lname{Hyp.ResidualBurnside} & $I[G]\twoheadrightarrow M_2(I)$ & tex proof of \lname{prop:s06-galois} &
\textbf{discharged} given absolute irreducibility ((H2))\\
\lname{Hyp.TamSpec} & \lname{lem:s06-tam}(a),(b) & Kim--Ota \cite[Def.~2.12]{KO}; \cite[Def.~3.3]{PW} & maths only\\
\lname{Hyp.HidaControlGL2} & control for $\T'$ at $P$, $\cO\otimes R\cong\T_k(N)_\m$ (GL$_2$ side; distinct from
\lname{Hyp.HidaControl}) & Hida \cite{Hi86b}; tex & cited\\
\lname{Hyp.KimOta} & \lname{thm:s06-KO} at one point & \cite[Thm.~1.1]{KO} & cited\\
\lname{Hyp.KimOtaRT} & \lname{prop:s06-levM} under (Lev) ($M=1$, or $M$ squarefree with $\bar\rho$ unramified at all $\ell\mid M$),
inclusion only; $N(\bar\rho)>1$ open & \cite[Thm.~3.14]{KO} & cited\\\bottomrule
\end{tabular}
\end{center}
Consumers: \lname{Setting.gor\_new} (\lname{HsiehGor}, \lname{LambdaJL}), \lname{Defect.galois}
(\lname{FreeRankTwo}, \lname{ResidualBurnside}), \lname{Defect.defect\_formula} (all of them, with
\lname{KimOta} $\lor$ \lname{KimOtaRT}), \lname{cor\_fixed}/\lname{cor\_allweights} (\lname{KimOta},
\lname{HidaControlGL2}); fixed-weight Gorenstein is derived by base change, not registered. Toy instances where feasible
(S06-L10b); they are consistency (non-vacuity) instances only, not evidence. \lname{Hyp.HsiehEta} will not be created. \lname{Hyp.LambdaJL} is the only existing hypothesis reused.

%======================================================================================================
\section{Inconsistencies found (and fixed)}\label{sec:hyp-incons}

Found while preparing the first version of this document by comparing \lname{HYPOTHESES.md}, the Lean files, the chapters,
the audits and the plans (state \lname{d214fb58}/\lname{cd0a268}). None affected the soundness of a Lean theorem: the type
of a theorem is what it is, and the axiom audit is clean. They affected the \emph{documentation} of what is assumed. All
eighteen were fixed by the work package S05-E1 (state of this document: after S05-E1); the list records what was wrong
and what was done.
\begin{enumerate}[leftmargin=1.6em,label=(I\arabic*),ref=I\arabic*]
\item\label{inc:invlim} \textbf{\lname{InvLimExact}.} The registry row lacked the connectivity premise (S04-L17c) and
  named ``compactness'' as a source; \lname{lim\_exact\_ribet(\_uncond)} take the literal \lname{hA : ∀ n, Surjective \dots}
  (the conclusion of \lname{Hyp.InvLimExact} given \lname{hconn}), not a term of that type. \emph{Fixed:} registry row
  (statement, source, ``used by'') and \lname{s04/lambdaribet.tex} now say exactly this; compactness is stated not to
  prove the surjectivity.
\item\label{inc:usedby} \textbf{Stale ``used by'' columns} (\lname{EichlerMass}, \lname{SA}, \lname{HS}, \lname{BTtree},
  \lname{PadicOrdDecomp}). \emph{Fixed} in the registry: the missing consumers were added (the four
  \lname{Ribet.*\_of\_SA\_HS} for \lname{SA}, \lname{HS}, \lname{BTtree}; \lname{sum\_inv\_wx\_Ohat}, \lname{mazur\_check},
  \lname{vp\_card\_XB\_Utower\_ge} for \lname{EichlerMass}; L10b, L17b and the \lname{\_uncond} twins for
  \lname{PadicOrdDecomp}); ``only'' removed from \lname{BTtree}.
\item\label{inc:planstaged} \textbf{Plan-only labels and \lname{plan/staged/} pointers.} \emph{Fixed:} \lname{\textbackslash label}s
  \lname{if:s02-SA} (\lname{s02/2-quotient.tex}), \lname{if:s02-tree} (\lname{s02/1-tree.tex}), \lname{if:s02-coh}
  (\lname{s02/3-cohomology.tex}) added; the seven Hyp docstrings (\lname{SA}, \lname{HS}, \lname{BTtree},
  \lname{SerreCoh}, \lname{SerreFree}, \lname{RamNrdSurj}, \lname{JLClassical}) point at the \lname{mathematics/} files.
\item\label{inc:nrd} \textbf{\lname{lem:s01-nrd}.} \emph{Fixed:} s01 now says ``\stproved{} on paper'' (the $\ell\mid D$ case rests
  on the cited local structure, Voight 13.3.11), drops the stale Lean pointer, and states the Lean status (formalized
  modulo \lname{Hyp.RamNrdSurj}); s02 says ``proved on paper in Lemma~s01-nrd, not formalized'', no longer \stcited.
\item\label{inc:mazur} \textbf{Mazur check for $q=2,3$.} \emph{Fixed:} \lname{s03/1-graphs.tex} says these cases are formalized
  through the $q\le3$ clause of \lname{Hyp.EllipticCount} (\lname{Ribet.mazur\_check}, every prime $q$); the registry row says so.
\item\label{inc:bt} \textbf{BTtree in stage 3.} \emph{Fixed:} \lname{s03/3-ribet.tex} l.~113 now reads
  $\HypSA\wedge\HypHS\wedge\mathrm{Hyp.BTtree}$ for the \lname{\_of\_SA\_HS} versions.
\item\label{inc:cr} \textbf{\lname{HsiehCR}.} \emph{Fixed:} \lname{\textbackslash lean} mark added in \lname{s04/hsieh.tex}; the
  registry consumer now reads: none in Lean, not taken by stage-6 (\lname{HsiehGor} has no CR premise).
\item\label{inc:mord} \textbf{\lname{MordLambdaFree}.} \emph{Fixed:} the status is ``tex proved'' in the Lean docstring,
  registry and chapter (the chapter proves (a)); consumer: none named in the stage-5/6 plans (grep), recorded as such; the
  R14 scope note is in the registry row.
\item\label{inc:brandt} \textbf{\lname{BrandtStabFree}.} \emph{Fixed:} registry gives one estimate (1000--1500 lines, Morita
  + PID, probed; the elementary route is noted as unprobed) and records the unused premise ``$D$ prime'' as an open design
  decision, and the weaker-than-lemma level restriction. Since then discharged (\lname{Hyp.brandtStabFree\_holds}).
\item\label{inc:ihara} \textbf{\lname{IharaH0}.} \emph{Fixed:} registry row says clause (b) was deferred to S04-L15a, that the
  operators now exist, and that (b) is still not stated. The Lean docstring of \lname{Hyp.IharaH0} says the same (S05-E1 review).
\item\label{inc:guard} \textbf{\lname{LambdaJL}.} \emph{Fixed:} \lname{thm:s01-lambdaJL} states the guard $\gcd(d,Mp)=1$.
\item \textbf{\lname{HidaControl}.} \emph{Fixed (documentation):} \lname{s04/control.tex} and the registry record the
  M9.3 sub-check: modules versus duals, the guard, the $U_1(M)$ restriction, and the unused \lname{hk}, \lname{hφ} (O1 is recorded,
  not applied in Lean).
\item \textbf{\lname{EllipticCount} row.} \emph{Fixed:} typo ``\lname{[D ≡ 2 mod 3\}}'' corrected in the registry.
\item \textbf{Import \lname{sorry}s.} \emph{Fixed:} \lname{SORRIES.md} now records the two import \lname{sorry}s, that they are
  outside the cone of \lname{Langlandscape}, and that two test files import one of them.
\item \textbf{Import instances not reachable.} \emph{Fixed:} the registry says the closed instances live in \lname{AlgJLTests},
  which is not a dependency; no instance in this project.
\item \textbf{Vendor bookkeeping.} \emph{Fixed before S05-E1} (commits \lname{6ae9b0c0}, \lname{0150bd23}): \lname{ATTRIBUTION.md}
  matches \lname{MANIFEST.tsv}; \lname{THIRD-PARTY-NOTICES.md} says Apache-2.0.
\item \textbf{Audit coverage.} \emph{Fixed:} \lname{Audit/PrintAxioms.lean} has lines for the six stage-2 Hyp definitions
  (\lname{SA}, \lname{HS}, \lname{BTtree}, \lname{SerreCoh}, \lname{SerreFree}, \lname{RamNrdSurj}) (to be run in the next
  TITAN round); the registry's ``4/4'' now names the three standard axioms.
\item \textbf{Bibliography.} \emph{Fixed:} Nikolov--Segal is in \lname{references.bib} (and removed from
  \lname{report/hypotheses.bib}); the duplicate Vign\'eras entry \lname{Vig80} is removed and the citations use \lname{Vigneras1980}.
\end{enumerate}

%======================================================================================================
\section{Exact Lean names}\label{sec:hyp-names}
\begingroup\raggedright

\paragraph{Registered hypotheses} (prefix \lname{Langlandscape.Hyp.} unless stated; files under
\lname{lean/Langlandscape/Langlandscape/Hyp/}):
\lname{GL2TF.Hyp.EichlerMass} (import; alias \lname{EichlerMass}), \lname{JLClassical}, \lname{LambdaJL} (with
\lname{GL2CuspData}, \lname{LambdaJLGen}), \lname{SA}, \lname{HS}, \lname{BTtree}, \lname{SerreCoh}, \lname{SerreFree},
\lname{RamNrdSurj}, \lname{IharaH0}, \lname{EllipticCount}, \lname{BrandtStabFree}, \lname{PadicOrdDecomp},
\lname{MordLambdaFree}, \lname{HidaControl}, \lname{HsiehCR}, \lname{InvLimExact}.

\paragraph{Discharge.} \lname{Langlandscape.Ordinary.padicOrdDecomp\_holds}, \lname{Langlandscape.Hyp.ramNrdSurj\_holds}
(with \lname{Hyp.ramNrdSurj\_of\_local}, \lname{Local.exists\_nrd\_eq\_unit\_of\_ne\_two}, \lname{Local.exists\_nrd\_eq\_unit\_two}).

\paragraph{Main consumers} (prefix \lname{Langlandscape.}):
\lname{Quat.card\_XB\_Utower}, \lname{Smoke.eichlerMass\_apply}, \lname{Ribet.sum\_inv\_wx\_Ohat},
\lname{Ribet.mazur\_check}, \lname{Ordinary.vp\_card\_XB\_Utower\_ge};
\lname{Quat.exists\_Ohat\_nrd\_eq}, \lname{Quat.nrdSurj\_\{Ohat,U0,U1,Utower\}},
\lname{Tree.isQAdmissible\_\{Ohat,U0,U1,Utower\}};
\lname{Tree.XBone\_toXB\_bijective}, \lname{Tree.BGConnected\_of\_surj}, \lname{Tree.BGConnected\_of\_SA\_HS},
\lname{Tree.surjType0\_of\_SA\_HS}, \lname{Tree.quotVEquiv\_of\_SA\_HS}, \lname{Tree.quotEEquiv\_of\_SA\_HS},
\lname{Tree.quotEquiv\_oO\_of\_SA\_HS}, \lname{Tree.quotEquiv\_tO\_of\_SA\_HS};
\lname{Ribet.ribet\_exact\_of\_SA\_HS}, \lname{Ribet.ribet\_rank\_of\_SA\_HS}, \lname{Ribet.ihara\_surj\_of\_SA\_HS},
\lname{Ribet.ihara\_ker\_const\_of\_SA\_HS};
\lname{Ribet.phiComb\_eisenstein\_Ohat}, \lname{Ribet.phiComb\_eisenstein\_DR}, \lname{Ribet.brandtZ\_dvd};
\lname{Ordinary.lim\_exact\_ribet}, \lname{Ordinary.lim\_exact\_ribet\_uncond} and the other \lname{\_uncond} theorems of
\lname{Ordinary/TowerUncond.lean}.

\paragraph{Tests that instantiate or link hypotheses} (prefix \lname{LanglandscapeTests.}):
\lname{S04.Hyp.hsiehCR\_toy}, \lname{S03.Mazur.test\_mazur\_check\_*}; files \lname{S02/SerreShape.lean},
\lname{S04/HypLam.lean}, \lname{S01/MassTower.lean}.

\par\endgroup

%======================================================================================================
\section{Refreshing this document}\label{sec:hyp-refresh}

This version is \emph{provisional} (2026-11-15, stage 4 closing, S04-L17c in integration). After the stage-4 final gate,
and again after stages 5 and 6:
\begin{enumerate}[leftmargin=1.5em]
\item Regenerate the status table from the repository root:
  \lname{python3 mathematics/report/tools/hyp\_table.py lean/Langlandscape/HYPOTHESES.md >
  mathematics/report/generated/hyp-table.tex}. The script warns on stderr about any hypothesis without a kind in its
  \lname{KIND} table; add new hypotheses (stage 5: \lname{X.char}, \lname{X.gal}, \lname{X.iso}; stage 6: nine names)
  to \lname{KIND} and \lname{STAGE} there, and remove the stage-5 names from its \lname{SKIP} set (S05-E1 skips
  them so that the stage-4 table stays as it is while they are still listed under the projected hypotheses).
\item Diff \lname{lean/Langlandscape/Langlandscape/Hyp/} and \lname{HYPOTHESES.md} against \lname{cd0a268}; update the
  affected entries of \S\ref{sec:hyp-entries}, the arrows of \lname{figures/hyp-deps.tex} and the events of
  \lname{figures/hyp-timeline.tex}. Move discharged hypotheses to \kinddis.
\item Replace the axiom-audit table of \S\ref{sec:hyp-intro} with the final gate's \lname{axcount.txt}; check
  \lname{pa.log} for \lname{sorryAx} (\lname{grep -c sorryAx}).
\item Re-check the items of \S\ref{sec:hyp-incons}; delete those that are fixed (all eighteen were fixed at S05-E1; new findings go here).
\item\sloppy Rebuild from \lname{mathematics/}:
  \lname{latexmk -pdf -interaction=nonstopmode -outdir=report report/standalone-hypotheses.tex}, then remove the aux files
  (\lname{latexmk -c -outdir=report report/standalone-hypotheses.tex}).
\end{enumerate}
